% Claim proof file for row claim_3_3 -- "AcyclicPreservedUnderDo".
% Auto-generated stub by scaffold/solve_chapter.py at row start. A manager
% agent dedicated to writing the TeX proof fills in the proof body below.
%
% This file is a sibling of `main.tex` inside `<subsection>/tex/`. The
% `subfiles` package lets it render both standalone and as part of
% `main.tex` via `\subfile{...}`.
%
% The statement is restated above the proof so the file is self-contained
% when rendered on its own. The proof must:
%   - Stay close to the lecture notes' proof if one exists (always search
%     the LN first for a `\begin{proof}` block immediately following the
%     claim; copy / lightly edit when present).
%   - Otherwise use the LN paradigm and cite prior claims by their `ref`.
%   - Be self-contained enough that a mathematician can follow it without
%     re-reading the LN.
%   - Have no `sorry`, `omitted`, or "obvious" shortcuts.

\documentclass[main]{subfiles}

\begin{document}
% ref: claim_3_3 (proof, with statement restated for readability)
% title: AcyclicPreservedUnderDo
\def\rowref{claim\_3\_3}
\phantomsection\label{claim_3_3}

% --- statement (restated) ---------------------------------------------------
\begin{Rem}[AcyclicPreservedUnderDo]
    Let $G = (J, V, E, L)$ be a CDMG (def \ref{def-cdmg}); throughout we use the notation of def \ref{not-cdmg} and recall the typing and axioms of def \ref{def-cdmg}, namely $J \cap V = \emptyset$, $E \subseteq (J \cup V) \times V$, $L \subseteq V \times V$, $L$ irreflexive ($(v_1, v_2) \in L \Rightarrow v_1 \neq v_2$), and $L$ symmetric ($(v_1, v_2) \in L \Leftrightarrow (v_2, v_1) \in L$). Let
    \[ W \subseteq J \cup V \]
    be a subset of nodes, and let $G_{\doit(W)} = (J_{\doit(W)}, V_{\doit(W)}, E_{\doit(W)}, L_{\doit(W)})$ be the hard-intervention CDMG of $G$ w.r.t.\ $W$, as constructed in def \ref{def:G_hard_intervention}. Both quantifiers ``for every CDMG $G$'' and ``for every $W \subseteq J \cup V$'' are read \emph{universally}: the side condition $W \subseteq J \cup V$ is the one inherited verbatim from def \ref{def:G_hard_intervention} (so $W$ is permitted to overlap the input-node set $J$), and the LN's free occurrence of $W$ in the claim body ranges over all such $W$ for which $G_{\doit(W)}$ is defined per that definition.

    \emph{Carrier matching.} Before stating the two conclusions, observe that the node carriers of $G$ and $G_{\doit(W)}$ coincide as subsets of the ambient universe:
    \[
        J_{\doit(W)} \cup V_{\doit(W)} \;=\; (J \cup W) \cup (V \sm W) \;=\; J \cup V,
    \]
    where the first equality is items i.\ and ii.\ of def \ref{def:G_hard_intervention} and the second equality uses $W \subseteq J \cup V$: the inclusion $(J \cup W) \cup (V \sm W) \subseteq J \cup V$ follows from $J \subseteq J \cup V$, $W \subseteq J \cup V$, and $V \sm W \subseteq V \subseteq J \cup V$; the reverse inclusion $J \cup V \subseteq (J \cup W) \cup (V \sm W)$ follows by case analysis on $x \in J \cup V$, namely $x \in J \Rightarrow x \in J \cup W$, $x \in V \cap W \Rightarrow x \in W \subseteq J \cup W$, and $x \in V \sm W \Rightarrow x \in V \sm W$. Consequently, a binary relation $<$ on $J \cup V$ is, without any re-indexing, the very same binary relation on $J_{\doit(W)} \cup V_{\doit(W)}$, and the predicate ``$<$ is a topological order'' (def \ref{def-topological-order}) is well-typed both for $G$ and for $G_{\doit(W)}$ with this single carrier.

    \emph{The remark.} If $G$ is acyclic (in the sense of def \ref{def-acylic}), then the following conjunction of sub-claims holds:
    \begin{enumerate}[label=(\alph*)]
        \item \emph{Acyclicity is preserved by $\doit(W)$.}
              The CDMG $G_{\doit(W)}$ is acyclic in the sense of def \ref{def-acylic}, i.e.\ for every $v \in J_{\doit(W)} \cup V_{\doit(W)} = J \cup V$, there does not exist any non-trivial directed walk from $v$ to itself in $G_{\doit(W)}$.
        \item \emph{Every topological order of $G$ is a topological order of $G_{\doit(W)}$.}
              For every strict total order $<$ on $J \cup V$,
              \[
                  \bigl(\, <\ \text{is a topological order of } G \,\bigr) \;\Longrightarrow\; \bigl(\, <\ \text{is a topological order of } G_{\doit(W)} \,\bigr),
              \]
              where both occurrences of ``topological order'' are in the sense of def \ref{def-topological-order} and the second occurrence uses the carrier-matching equality $J_{\doit(W)} \cup V_{\doit(W)} = J \cup V$ above to view $<$ unchanged as a strict total order on the carrier of $G_{\doit(W)}$. The LN's indefinite article in ``a topological order for $G$ is also one for $G_{\doit(W)}$'' is read \emph{universally} (every such $<$, not merely some such $<$).
    \end{enumerate}
    The bidirected-edge set $L$ plays no role in either sub-claim: acyclicity (def \ref{def-acylic}) is defined via the non-existence of non-trivial directed walks (which use only $E$), and a topological order (def \ref{def-topological-order}) is defined via the parent-set $\Pa^G$ (def \ref{def_3_5}, item~i., which is defined from $E$ alone).
\end{Rem}

% --- proof ------------------------------------------------------------------
\begin{proof}
The proof rests on a single observation about directed edges, established as a lemma below; both sub-claims then follow by short structural arguments. Throughout the proof we use the carrier-matching equality $J_{\doit(W)} \cup V_{\doit(W)} = J \cup V$ already established in the statement block above, so that the directed-edge set $E_{\doit(W)}$ and the parent-set $\Pa^{G_{\doit(W)}}$ are both well-typed on the same node carrier $J \cup V$ as their $G$-counterparts.

\medskip\noindent\emph{(0) Edge-inclusion lemma: $E_{\doit(W)} \subseteq E$.}
By item iii.\ of def \ref{def:G_hard_intervention}, the directed-edge set of $G_{\doit(W)}$ is constructed as a set difference of $E$:
\[
    E_{\doit(W)} \;:=\; E \sm \lC v \tuh w \,\middle|\, v \in G,\ w \in W \rC,
\]
i.e.\ every element of $E_{\doit(W)}$ is by definition an element of $E$ from which a certain subset has been removed. Hence
\[
    E_{\doit(W)} \;\subseteq\; E \qquad \text{as subsets of } (J \cup V) \times V.
\]
Consequently, if
\[
    p \;=\; (v_0, a_0, v_1, a_1, \dots, a_{n-1}, v_n)
\]
is a directed walk in $G_{\doit(W)}$ in the sense of def \ref{def:walks}, item~ii.\ -- that is, $v_0, v_1, \dots, v_n \in J_{\doit(W)} \cup V_{\doit(W)}$ and $a_k = (v_k, v_{k+1}) \in E_{\doit(W)}$ for every $k \in \{0, 1, \dots, n - 1\}$ -- then \emph{the same tuple} $p$ is also a directed walk in $G$: by the carrier-matching equality each $v_k$ lies in $J_{\doit(W)} \cup V_{\doit(W)} = J \cup V$, so the vertex-typing requirement of def \ref{def:walks}, item~ii.\ for $G$ is met; and by the inclusion $E_{\doit(W)} \subseteq E$ just derived, each edge $a_k = (v_k, v_{k+1}) \in E_{\doit(W)} \subseteq E$, so the edge-typing requirement of def \ref{def:walks}, item~ii.\ for $G$ is met. The lift is verbatim: same vertex sequence $(v_0, v_1, \dots, v_n)$, same edge sequence $(a_0, a_1, \dots, a_{n-1})$, same length $n$, no re-indexing of any kind. We make no assertion about the bidirected-edge set $L_{\doit(W)}$ vis-\`a-vis $L$ here; neither sub-claim requires it, because acyclicity (def \ref{def-acylic}) is defined via directed walks (whose edges are all in $E$) and the parent-set $\Pa^G$ (def \ref{def_3_5}, item~i.) is defined from $E$ alone.

\medskip\noindent\emph{(a) Acyclicity is preserved by $\doit(W)$.}
Assume $G$ is acyclic in the sense of def \ref{def-acylic}: for every $v \in J \cup V$, there does not exist any non-trivial directed walk from $v$ to itself in $G$. We must show the analogous statement for $G_{\doit(W)}$ on its carrier $J_{\doit(W)} \cup V_{\doit(W)}$. By the carrier-matching equality $J_{\doit(W)} \cup V_{\doit(W)} = J \cup V$ from the statement above, this carrier is literally the same set $J \cup V$, so we proceed by contradiction. Suppose, for contradiction, that there exist some $v \in J_{\doit(W)} \cup V_{\doit(W)} = J \cup V$ and a non-trivial directed walk
\[
    p \;=\; (v_0, a_0, v_1, a_1, \dots, a_{n-1}, v_n) \qquad \text{in } G_{\doit(W)},
\]
with $n \ge 1$ (non-trivial), $v_0 = v_n = v$, and each $a_k = (v_k, v_{k+1}) \in E_{\doit(W)}$. By the edge-inclusion lemma~(0), the identical tuple $p$ is also a directed walk in $G$, with the very same length $n \ge 1$ and the very same end-nodes $v_0 = v_n = v \in J \cup V$. This is therefore a non-trivial directed walk from $v$ to itself in $G$, contradicting the acyclicity of $G$. Hence no such walk $p$ exists in $G_{\doit(W)}$, and $G_{\doit(W)}$ is acyclic in the sense of def \ref{def-acylic}.

\medskip\noindent\emph{(b) Every topological order of $G$ is a topological order of $G_{\doit(W)}$.}
Let $<$ be any strict total order on $J \cup V$ that is a topological order of $G$ in the sense of def \ref{def-topological-order}; concretely, $<$ is a binary relation on $J \cup V$ that is irreflexive, transitive, and satisfies trichotomy ($v < w$, $v = w$, or $w < v$ for every $v, w \in J \cup V$), and satisfies the parent-precedence clause
\[
    \forall\, v, w \in J \cup V :\quad v \in \Pa^G(w) \;\Longrightarrow\; v < w.
\]
We verify, in turn, that $<$ is a strict total order on $J_{\doit(W)} \cup V_{\doit(W)}$ and that the parent-precedence clause for $G_{\doit(W)}$ holds for $<$.

\smallskip\noindent\emph{Strict total order on the carrier of $G_{\doit(W)}$.}
By the carrier-matching equality $J_{\doit(W)} \cup V_{\doit(W)} = J \cup V$ established in the statement block, the carrier of $G_{\doit(W)}$ as a set is identically the carrier of $G$. Hence $<$ -- viewed unchanged as the very same binary relation -- is already a binary relation on $J_{\doit(W)} \cup V_{\doit(W)}$ (no domain restriction or extension is required). Irreflexivity, transitivity, and trichotomy of $<$ on $J \cup V$ transfer to $<$ on $J_{\doit(W)} \cup V_{\doit(W)}$ because the universally-quantified variables in each of those three axioms range over literally the same set $J \cup V = J_{\doit(W)} \cup V_{\doit(W)}$. Hence $<$ is a strict total order on $J_{\doit(W)} \cup V_{\doit(W)}$.

\smallskip\noindent\emph{Parent-set inclusion $\Pa^{G_{\doit(W)}}(w) \subseteq \Pa^G(w)$ for every $w \in J \cup V$.}
Fix $w \in J_{\doit(W)} \cup V_{\doit(W)} = J \cup V$. Applying the set-builder definition of the parent set (def \ref{def_3_5}, item~i.) to the CDMG $G_{\doit(W)}$, and using the carrier equality above,
\[
    \Pa^{G_{\doit(W)}}(w) \;=\; \lC v \in J_{\doit(W)} \cup V_{\doit(W)} \,\middle|\, (v, w) \in E_{\doit(W)} \rC
                          \;=\; \lC v \in J \cup V \,\middle|\, (v, w) \in E_{\doit(W)} \rC.
\]
By the edge-inclusion lemma~(0), the membership condition $(v, w) \in E_{\doit(W)}$ implies $(v, w) \in E$. Hence
\[
    \lC v \in J \cup V \,\middle|\, (v, w) \in E_{\doit(W)} \rC
    \;\subseteq\;
    \lC v \in J \cup V \,\middle|\, (v, w) \in E \rC
    \;=\;
    \Pa^G(w),
\]
where the last equality is def \ref{def_3_5}, item~i., applied to the CDMG $G$. Combining the two displayed identities gives the desired inclusion $\Pa^{G_{\doit(W)}}(w) \subseteq \Pa^G(w)$.

\smallskip\noindent\emph{Parent-precedence on $G_{\doit(W)}$.}
Let $v, w \in J_{\doit(W)} \cup V_{\doit(W)} = J \cup V$ and suppose $v \in \Pa^{G_{\doit(W)}}(w)$. By the parent-set inclusion just derived, $v \in \Pa^G(w)$. Since $<$ is a topological order of $G$, the parent-precedence clause for $G$ yields $v < w$. Thus
\[
    \forall\, v, w \in J_{\doit(W)} \cup V_{\doit(W)} :\quad v \in \Pa^{G_{\doit(W)}}(w) \;\Longrightarrow\; v < w,
\]
which is exactly the parent-precedence clause of def \ref{def-topological-order} for $G_{\doit(W)}$.

\smallskip
Combining the strict-total-order check and the parent-precedence check, $<$ satisfies both conjuncts of def \ref{def-topological-order} for $G_{\doit(W)}$ and is therefore a topological order of $G_{\doit(W)}$. Since $<$ was an arbitrary topological order of $G$, every topological order of $G$ is a topological order of $G_{\doit(W)}$, establishing sub-claim (b).

\medskip\noindent\emph{Closing remark on the role of the acyclicity hypothesis.}
The argument for sub-claim (b) used only the edge-inclusion lemma~(0) and the carrier-matching equality; the acyclicity hypothesis on $G$ was never invoked. This is structural: the parent set $\Pa^G$ (def \ref{def_3_5}, item~i.) is defined from the directed-edge set $E$ via the membership condition $(v, w) \in E$, so removing directed edges (as $\doit(W)$ does in item iii.\ of def \ref{def:G_hard_intervention}) can only shrink the parent set, which is the only force needed to transfer parent-precedence. Hence sub-claim (b) holds for every CDMG $G$ and every $W \subseteq J \cup V$, with or without acyclicity. The LN's packaging of (a) and (b) under the common antecedent ``If $G$ is acyclic'' is preserved verbatim in the statement and proof above; the antecedent is genuinely used only in sub-claim (a) (where the lifted self-walk in $G$ is required to contradict $G$'s acyclicity).
\end{proof}

\end{document}
